\documentclass[aspectratio=169, 11pt]{beamer}

% --- Packages & Setup ---
\usepackage[utf8]{inputenc}
\usepackage[T1]{fontenc}
\usepackage{xcolor}
\usepackage{booktabs} 
\usepackage{graphicx}
\usepackage{tikz} % Required for the corner box
\usetikzlibrary{calc}

\graphicspath{{../}{./}}

% Try loading Open Sans; fallback if missing
\IfFileExists{opensans.sty}{\usepackage[default,scale=0.95]{opensans}}{}

% --- Branding Colors ---
\definecolor{amritaMaroon}{HTML}{A4123F} 

% --- Theme Setup ---
\usetheme{Madrid}
\usecolortheme[named=amritaMaroon]{structure}
\setbeamertemplate{navigation symbols}{}
\setbeamertemplate{footline}[frame number]

\begin{document}

% --- CUSTOM TITLE SLIDE (Research) ---
\begin{frame}
    % Corner Box Indicator
    \begin{tikzpicture}[remember picture,overlay]
        \node[fill=amritaMaroon, text=white, font=\bfseries\footnotesize, anchor=north east, rounded corners=2pt, inner sep=5pt] 
        at ($(current page.north east) + (-0.2cm,-0.2cm)$) 
        {Type 1: Research Project};
    \end{tikzpicture}

    \centering
    \includegraphics[width=2.8cm]{logo.png} 
    \vspace{0.1cm}
    
    {\large \textbf{\textcolor{amritaMaroon}{Something TDCCEA}}} \\
    \vspace{0.05cm}
    {\footnotesize \textbf{23CSE498 -- Project Phase 2 (Panel Review 1)} \quad | \quad \textbf{Team Number: 99}} \\ 
    \vspace{0.15cm}

    % Student Table
    \begin{table}
        \centering
        \footnotesize 
        \renewcommand{\arraystretch}{1.05}
        \begin{tabular}{|c|c|l|}
            \hline
            \textbf{S.No} & \textbf{Register Number} & \textbf{Student Name} \\
            \hline
            1 & CB.SC.U4CSE23215 & Darun R G \\
            \hline
            2 & CB.SC.U4CSE23232 & Pradeep Sasidharan \\
            \hline
            3 & CB.SC.U4CSE23235 & R Anusanth \\
            \hline
            4 & CB.SC.U4CSE23256 & Sathya Roopan M \\
            \hline
            5 & CB.SC.U4CSE23152 & Vinaayak Kanagaraj \\
            \hline
        \end{tabular}
    \end{table}
    \vspace{0.1cm}
    
    {\footnotesize Guide: \textbf{Dr. G Jeyakumar}, Professor, Department of CSE} \\
    \vspace{0.1cm}
    {\scriptsize \today}
\end{frame}

% --- SLIDE 1: METHODOLOGY & SOLUTION DESIGN (Rubric: 5 Marks) ---




\begin{frame}{Agenda}\tableofcontents\end{frame}




\section{Problem Statement \& Research Gaps}
% ════════════════════════════════════════════════════════════

\begin{frame}{Problem Statement}
\begin{block}{Problem}
Strategy Optimization and Co-Evolution of Competitive Game Agents Using Evolutionary Algorithms.
\end{block}
\vspace{0.15cm}
\small
\textbf{Abstract:} We introduce a game system where agents learn and improve continuously rather than being trained once. Agents are evaluated using \textbf{weighted performance histories} across multiple generations --- capturing adaptive capability over time.

\vspace{0.15cm}
Instead of adapting to opponent behavior alone, agents evolve in response to shifts in the \textbf{game environment itself} --- changes in map configuration, resource distribution, and competitive conditions.

\vspace{0.15cm}
A \textbf{diversity-preservation mechanism} prevents populations from converging to similar behaviors, while a \textbf{Hall-of-Fame archive} ensures agents are tested against historically strong opponents.

\vspace{0.15cm}
Together, these mechanisms produce game agents that remain competitive, strategically diverse, and stable as the environment evolves --- validated on the microRTS simulation environment.
\end{frame}

\begin{frame}{Research Gaps --- Phase 1 (Broad)}
\small
\begin{itemize}
    \item Most existing approaches are \textbf{game-specific}, limiting generalization across different environments.
    \item Many methods rely on \textbf{manual tuning}, hand-crafted rules, or fixed parameter settings, reducing scalability.
    \item Current approaches often optimize \textbf{single objectives} (e.g., win rate) without considering long-term adaptation or evolving opponents.
    \item Limited focus on \textbf{time-dependent performance} and strategy evolution across generations.
    \item Need for more \textbf{adaptive, trajectory-aware, and coevolutionary frameworks} for dynamic game environments.
\end{itemize}
\vspace{0.2cm}

\end{frame}

\begin{frame}{Research Gaps --- Phase 2}
\scriptsize
\begin{itemize}
    \item \textbf{\textcolor{amritaMaroon}{Gap 1: Red Queen Dynamics \& Loss of Gradients}}\\
    Both populations adapt but neither improves in absolute terms --- each merely keeps pace. (Rawal et al.\ 2010; Elfeky et al.\ 2022)
    \vspace{0.15cm}
    
    \item \textbf{\textcolor{amritaMaroon}{Gap 2: Cycling, Forgetting \& Lack of Diversity}}\\
    Co-evolution suffers from forgetting previously discovered solutions and cycling between strategies. (Runarsson \& Lucas 2005; Hevia Fajardo et al.\ 2023)
    \vspace{0.15cm}
    
    \item \textbf{\textcolor{amritaMaroon}{Gap 3: Defensive Convergence \& Strategy Stagnation}}\\
    Agents converge on safe defensive strategies without fitness shaping. Unbalanced games cause rapid stagnation on a single dominating strategy. (Dziuk \& Miikkulainen 2011; Leigh et al.\ 2008)
    \vspace{0.15cm}
    
    \item \textbf{\textcolor{amritaMaroon}{Gap 4: Difficulty of Measuring True Fitness Progress}}\\
    Fitness values reflect relative performance against the current opponent, not absolute quality --- apparent improvements may be opponent drift. (Lucas 2006; Hevia Fajardo et al.\ 2024)
    \vspace{0.15cm}
    
    \item \textbf{\textcolor{amritaMaroon}{Gap 5: Self-Adaptive Mutation Without Environmental Context}}\\
    Adaptive mechanisms often respond only to pairwise comparisons, ignoring the broader environmental context of how the competitive landscape has shifted. (Hevia Fajardo et al.\ 2024)
\end{itemize}
\end{frame}



% --- SLIDE 1: METHODOLOGY & SOLUTION DESIGN (Rubric: 5 Marks) ---
\section{Methodology \& Solution Design}
\begin{frame}{1. Methodology \& Solution Design --- Part A}
    \framesubtitle{Rubric: Workflow, Architecture, Algorithms/Mathematical Model \& Evaluation Metrics (5 Marks)}
    
    \begin{itemize}
        \item \textbf{Theoretical Foundation:} Trajectory-Driven Competitive Co-Evolution (TD-CCEA) --- fitness is scored against a \textbf{weighted trajectory} of past environment states $c_t=\frac{1}{N}\sum_i\theta_i^t$, not just the current opponent, giving adaptive mutation $\mu_i = \mu_{base}(1-e_i^t)+\epsilon$. This reduces Red Queen dynamics and opponent-drift bias.
        
        \item \textbf{Core Algorithms:} Four variants compared:
        \begin{itemize}
            \item TD-CCEA (base): trajectory memory + adaptive mutation
            \item SA-TDCCEA: sensitivity-adaptive mutation (per-gene bounds $\mu_{min}{=}0.08$, $\mu_{max}{=}0.20$)
            \item CCE-TDCCEA: 8D macro (economy) / 12D micro (combat) cooperative decomposition
            \item Standard CCEA: no trajectory memory, fixed mutation $\mu{=}0.15$ (baseline)
        \end{itemize}
        
        
    \end{itemize}
\end{frame}
 
\begin{frame}{1. Methodology \& Solution Design --- Part B}
    \framesubtitle{Rubric: Workflow, Architecture, Algorithms/Mathematical Model \& Evaluation Metrics (5 Marks)}
    
    \begin{itemize}
        
        
        \item \textbf{Environment \& Preprocessing:} microRTS simulation engine (Java-based strategy game). Dynamic 48$\times$48 maps:
        \begin{itemize}
            \item Progressive blockades: 6--14 walls added each generation (early stages $<$ mid $<$ late)
            \item Mixed wall types: 60\% unbreakable terrain + 40\% destructible (HP=5)
            \item BFS-validated: every wall placement checked for base-to-base path reachability
            \item \Opponent rotation: 7 built-in bots used for validation 
        \end{itemize}
    \end{itemize}
\end{frame}
 
\begin{frame}{1. Methodology \& Solution Design --- Part C}
    \framesubtitle{Rubric: Workflow, Architecture, Algorithms/Mathematical Model \& Evaluation Metrics (5 Marks)}
    \footnotesize
    \begin{columns}[t]
        \column{0.5\textwidth}
        \textbf{Evaluation Metrics:}
        \begin{itemize}\footnotesize
            \item \textbf{vs.\ 7 Built-in Bots}: Win rate (\%) per bot; overall average across all bots
            \item \textbf{Head-to-Head Tournament}: Round-robin (each pair plays 10 games with alternating sides); W/L/D record per matchup; aggregate win rate
            \item \textbf{Statistical Significance}: Wilcoxon signed-rank test on per-bot win rates (CCE-TDCCEA vs.\ each baseline); $p{<}0.05$ threshold
            \item \textbf{Ablation Analysis}: Contribution of trajectory memory, cooperative decomposition, and sensitivity-adaptive mutation measured separately
        \end{itemize}
        
        \textbf{Key Hyperparameters:}
        \begin{itemize}\footnotesize
            \item $N{=}20$ (population size)
            \item $T{=}30$ (generations per coevolution run)
            \item $\mu_{\text{base}}{=}0.15$, $\epsilon{=}0.03$ (adaptive mutation)
            \item $d_{\max}{=}\sqrt{20}$ (Euclidean distance norm, 20D weight space)
        \end{itemize}
 
        \column{0.5\textwidth}
        \centering
        \fbox{\begin{minipage}{6.8cm}
            \centering
            \vspace{0.8cm}
            \textbf{[Insert Methodology / Algorithmic Workflow Diagram]} \\
            \scriptsize (Show: Input Data $\to$ Feature Extraction $\to$ Model $\to$ Evaluation)
            \vspace{0.8cm}
        \end{minipage}}
    \end{columns}
\end{frame}
 

% --- SLIDE 2: IMPLEMENTATION PROGRESS (~60%) (Rubric: 8 Marks) ---
\section{Implementation Progress}
\begin{frame}{2. Implementation Progress ($\approx$60\%) \& Module Functionality}
    \framesubtitle{Rubric: Core Modules Implemented, Functional \& Integrated (8 Marks)}
    
    {\small
    \begin{itemize}
        \item \textbf{Implementation Status Overview ($\approx$60\% Completed):}
        
        {\scriptsize
        \begin{table}
            \centering
            \renewcommand{\arraystretch}{1.05}
            \begin{tabular}{p{3.5cm} p{4.5cm} c}
                \toprule
                \textbf{Module / Component} & \textbf{Current Implementation State} & \textbf{Status} \\
                \midrule
                Core TD-CCEA Algorithm & Trajectory scoring $e_i^t$, adaptive mutation $\mu_i$, diversity-preserving pruning & Functional \\
                SA-TDCCEA Variant & Offline sensitivity profiling $\to$ per-gene mutation bounds ($\mu_{min}{=}0.08$, $\mu_{max}{=}0.20$) & Functional \\
                CCE-TDCCEA Variant & 8D macro / 12D micro cooperative sub-population decomposition & In Progress \\
                Dynamic Environment (BFS) & Progressive blockade generation with BFS-validated base-to-base reachability & Functional \\
                microRTS Integration \& Game Server & EVOLVED-vs-BOT and EVOLVED-vs-EVOLVED batch game runner (Java $\leftrightarrow$ Python) & Functional \\
                Benchmark \& Head-to-Head Harness & Parallel game pool, round-robin tournament scaffolding vs. Standard CCEA \& classical EAs & In Progress \\
                \bottomrule
            \end{tabular}
        \end{table}
        }
        
        \item \textbf{Core Functionality Demonstration:} 100\% win rate validated for core TD-CCEA on static 8$\times$8 and 16$\times$16 microRTS maps; dynamic 32$\times$32 blockade environment and cross-algorithm coevolution runs are executing on GitHub-tracked commits.
    \end{itemize}
    }
    
\end{frame} 


% ════════════════════════════════════════════════════════════





% --- SLIDE 3: TECHNICAL KNOWLEDGE & DESIGN CHOICES (Rubric: 5 Marks) ---
\section{Technical Knowledge}
\begin{frame}{3. Technical Knowledge \& Design Decisions}
    \framesubtitle{Rubric: Understanding of Methodology, Implementation \& Design Choices (5 Marks)}
    
    \scriptsize
    \setlength{\itemsep}{1.5pt}
    \setlength{\parskip}{0pt}
    
    \begin{itemize}
        \item \textbf{Algorithmic Justification:} Standard co-evolutionary EAs suffer from Red Queen dynamics, forgetting, and defensive strategy stagnation (Gaps 1--3) because fitness is judged only against the \emph{current} opponent. TD-CCEA instead scores agents against a \textbf{weighted trajectory of past environment states} $c_t = \frac{1}{N}\sum_i \theta_i^t$, giving a fitness signal that reflects sustained performance rather than opponent drift.

        \item \textbf{Mathematical Formulation:}
        \[
        e_i^t = \sum_k \hat{\alpha}_k f_i^{(k)},\quad
        \hat{\alpha}_k = \frac{\alpha_k}{\sum_j \alpha_j},\quad
        \alpha_k = 1-\frac{d(c_t,c_k)}{d_{max}},
        \qquad
        \mu_i = \mu_{base}(1-e_i^t)+\epsilon
        \]
        \vspace{-6pt}
        
        High trajectory score $e_i^t \Rightarrow$ low mutation (exploit);
        low $e_i^t \Rightarrow$ high mutation (explore).

        \item \textbf{Design Choices \& Trade-offs:}
        \begin{itemize}
            \scriptsize
            \setlength{\itemsep}{1pt}
            \item Pruning thresholds $\delta_f$ (weak agents) and $\delta_s$ (redundant agents) trade population diversity against convergence speed.
            \item Hall-of-Fame archive capped at $A_{max}=15$ --- large enough to retain historically strong strategies, small enough to keep evaluation cost tractable.
            \item CCE-TDCCEA's 8D/12D macro-micro split trades some cross-weight coordination for faster, more specialized sub-population search.
            \item SA-TDCCEA's per-gene mutation bounds ($\mu_{min}=0.08$, $\mu_{max}=0.20$) trade a one-time offline sensitivity-profiling cost for finer-grained exploration control than a single global $\mu$.
        \end{itemize}

        \item \textbf{Preliminary Observations:} %Core TD-CCEA reaches 100\% win rate against baseline microRTS bots on static 8$\times$8 / 16$\times$16 maps, confirming the trajectory-scoring and adaptive-mutation mechanisms behave as intended before scaling up to the dynamic 32$\times$32 blockade environment.
    \end{itemize}
\end{frame}





































% ════════════════════════════════════════════════════════════

% ════════════════════════════════════════════════════════════
\section{Core Algorithm Recap}
% ════════════════════════════════════════════════════════════

\begin{frame}{Core Algorithm --- TD-CCEA}
\small
\textbf{Environment meta-state} (competitive landscape fingerprint):
\[ c_t = \frac{1}{N}\sum_{i=1}^{N} \theta_i^t \]
\textbf{Trajectory score} (smooths noisy fitness):
\[ e_i^t = \sum_k \hat{\alpha}_k \cdot f_i^{(k)},\quad
   \hat{\alpha}_k = \frac{\alpha_k}{\sum_j \alpha_j},\quad
   \alpha_k = 1 - \frac{d(c_t, c_k)}{d_{max}} \]
\textbf{Adaptive mutation} (dynamic --- not static):
\[ \mu_i = \mu_{base}\cdot(1 - e_i^t) + \epsilon \]
\begin{itemize}
    \item High fitness $\rightarrow$ low mutation (exploit). Low fitness $\rightarrow$ high mutation (explore).
    \item \textbf{Pruning:} weak ($\bar{e}_i<\delta_f$) and redundant ($S(i,j)>\delta_s$) agents removed each generation.
    \item \textbf{Hall-of-Fame Archive:} top-3 per generation, bounded to $A_{max}=15$.
\end{itemize}
\end{frame}

\begin{frame}{Algorithm Representation}
\begin{block}{TD-CCEA}
Full iterative workflow: fitness evaluation with trajectory buffering, expressibility scoring, environment-aware adaptation, diversity control, archive maintenance, and reproduction.
\begin{center}
    \includegraphics[width=0.55\linewidth]{Algo.png}
\end{center}
\end{block}
\end{frame}

% ════════════════════════════════════════════════════════════
\section{Algorithm Variants}
% ════════════════════════════════════════════════════════════

\begin{frame}{Variant 1: SA-TDCCEA (Sensitivity-Adaptive)}
\begin{block}{Key Idea}
Not all genes are equally important. SA-TDCCEA uses sensitivity analysis to assign \textbf{per-gene mutation rates}.
\end{block}
\vspace{0.15cm}
\begin{itemize}
    \item \textbf{Sensitivity profiling:} For each weight $w_j$, measure $\Delta W_j = |\text{win}(w_j{=}1) - \text{win}(w_j{=}0)|$
    \item High-influence genes $\rightarrow$ small mutation ($\mu_{min}=0.08$)
    \item Low-influence genes $\rightarrow$ large mutation ($\mu_{max}=0.20$)
    \item Still scaled by trajectory score $e_i^t$
\end{itemize}
\vspace{0.15cm}

\end{frame}

\begin{frame}{Variant 2: CCE-TDCCEA (Cooperative Co-Evolution)}
\begin{block}{Key Idea}
Decompose the 20D strategy vector into two cooperating sub-populations evolved independently but evaluated jointly.
\end{block}
\vspace{0.15cm}
\begin{columns}[t]
    \column{0.48\textwidth}
    \textbf{Population A: Macro/Economy (8D)}
    \begin{itemize}\small
        \item worker\_ratio, max\_workers
        \item expand\_threshold, barracks\_timing
        \item economy\_pressure, scout\_priority
        \item second\_barracks, worker\_fight\_ratio
    \end{itemize}
    \column{0.48\textwidth}
    \textbf{Population B: Micro/Combat (12D)}
    \begin{itemize}\small
        \item light/heavy/ranged preference
        \item aggression, attack\_threshold
        \item defend\_radius, target\_priority
        \item defense\_trigger, counter\_rush
        \item tech\_switch, group\_attack, focus\_fire
    \end{itemize}
\end{columns}
\vspace{0.15cm}
\textbf{Cooperative Handshake:} Each paired with top-3 from other sub-population.\\

\end{frame}





% --- SLIDE 4: TECHNICAL ENRICHMENT ACTIVITY CHOICES ---
\section{Technical Enrichment}
\begin{frame}{4. Technical Enrichment Activity Choices}
    \framesubtitle{Professional Online Certification Details (Individual Mandatory: 10--20 Hours)}
    \begin{itemize}
        \item \textbf{Individual MOOC / Online Certification Selection:}
    \end{itemize}
    \vspace{-0.1cm}
    \begin{table}
        \centering
        \footnotesize
        \renewcommand{\arraystretch}{1.15}
        \begin{tabular}{|c|p{2.8cm}|p{5.0cm}|p{2.3cm}|c|}
            \hline
            \textbf{S.No} & \textbf{Student Name} & \textbf{Chosen MOOC / Course Title} & \textbf{Platform} & \textbf{Hours} \\
            \hline
            1 & Darun R G & [Course Title relevant to Project] & Coursera & 15 Hrs \\
            \hline
            2 & Pradeep Sasidharan & [Course Title relevant to Project] & NPTEL & 20 Hrs \\
            \hline
            3 & R Anusanth & [Course Title relevant to Project] & edX / Coursera & 12 Hrs \\
            \hline
            4 & Sathya Roopan M & [Course Title relevant to Project] & AWS / Google & 16 Hrs \\
            \hline
            5 & Vinaayak Kanagaraj & [Course Title relevant to Project] & Coursera / NPTEL & 15 Hrs \\
            \hline
        \end{tabular}
    \end{table}
    \vspace{0.15cm}
    \begin{itemize}
        \item {\scriptsize \textbf{Requirement:} Each student must complete one approved 10--20 hour online certification from a recognized platform (Coursera, NPTEL, edX, etc.) relevant to the project.}
    \end{itemize}
\end{frame}









% ──────────────────────────────────────────────────────────────
%  SECTION: Phase II — Dynamic Environment & Coevolution
% ──────────────────────────────────────────────────────────────
\section{Phase II: Dynamic Environment \& Coevolution}
% FIXED: 32×32 → 48×48, 3–7 → 6–14 walls, removed "Shifting Economics"
\begin{frame}{Dynamic Environment --- Why It Matters}
\begin{block}{The Static Environment Problem}
A static map allows agents to win by memorizing one fixed strategy, which doesn't test true adaptation. In real games, the terrain changes every match. By progressively adding blockades each generation, we force agents to continuously adapt rather than memorize
\end{block}
\vspace{0.15cm}
\textbf{Phase II Extension --- continuously shifting environment:}
\begin{enumerate}
    \item \textbf{Progressive Blockades:} Each generation adds 6--14 walls to the 48$\times$48 map
    \item \textbf{Mixed Wall Types:} 60\% unbreakable terrain + 40\% destructible (HP=5)
    \item \textbf{BFS Validation:} Every wall placement verified for base-to-base connectivity
\end{enumerate}
\vspace{0.1cm}
\end{frame}
% FIXED: walls 3/5/7 → 6/10/14, Gen 15–30 → Gen 15–22, 168 → 224


\begin{frame}{Progressive Blockade Generation}
    \centering
    
    \vspace{0.15cm}

    \resizebox{0.98\textwidth}{!}{%
    \begin{tikzpicture}[
        box/.style={
            rectangle,
            rounded corners=8pt,
            draw=amritaMaroon,
            fill=amritaMaroon!8,
            minimum height=1.35cm,
            text width=3.0cm,
            align=center,
            font=\small
        },
        arrow/.style={
            -{Latex[length=2.5mm]},
            thick,
            amritaMaroon
        }
    ]

        % Four progressive stages
        \node[box] (early) at (0,0) {
            \textbf{Early}\\[-1pt]
            Gen 1--10\\[-1pt]
            +6 walls/gen\\[-1pt]
        };

        \node[box] (mid) at (3.65,0) {
            \textbf{Mid}\\[-1pt]
            Gen 11--20\\[-1pt]
            +10 walls/gen\\[-1pt]
        };

        \node[box] (late) at (7.30,0) {
            \textbf{Late}\\[-1pt]
            Gen 21--30\\[-1pt]
            +14 walls/gen\\[-1pt]
        };

        \node[box] (final) at (10.95,0) {
            \textbf{Final State}\\[-1pt]
            $\sim$224 total walls\\[-1pt]
            Complex battlefield
        };

        % Arrows
        \draw[arrow] (early.east) -- (mid.west);
        \draw[arrow] (mid.east) -- (late.west);
        \draw[arrow] (late.east) -- (final.west);

    \end{tikzpicture}%
    }

    \vspace{0.35cm}

    \begin{block}{BFS Reachability Guarantee}
        \centering
        \small
        Every candidate wall is validated before placement ---
        no soft-locked game states.
        
        \vspace{-0.05cm}
        
        \[
        \mathrm{BFS\_Reachable}
        \left(
        \mathcal{W}_{\mathrm{all}}\cup\{(x,y)\},
        \mathrm{Base}_0,
        \mathrm{Base}_1
        \right)
        =
        \mathrm{True}
        \]
    \end{block}

\end{frame}





% FIXED: 2 algorithms → 4 algorithms, pop 15 → 20, W=5 → W=7
\begin{frame}{Competitive Coevolution --- Arms Race}
\begin{block}{Four Algorithms Co-Evolve Simultaneously}
All four algorithms evolve against \textbf{each other's current best agent}. Each generation, every algorithm faces the best weights of the other three as opponents.
\end{block}
\vspace{0.15cm}\small
\begin{columns}[t]
    \column{0.48\textwidth}
    \textbf{TD-CCEA / CCE-TDCCEA / SA-TDCCEA}
    \begin{itemize}
        \item Population of 20 agents
        \item Trajectory memory ($W=7$)
        \item Adaptive mutation
        \item Diversity-preserving pruning
        \item Hall-of-Fame archive
    \end{itemize}
    \column{0.48\textwidth}
    \textbf{Standard CCEA (Baseline)}
    \begin{itemize}
        \item Population of 20 agents
        \item No trajectory memory (memoryless)
        \item Fixed Gaussian mutation ($\mu=0.15$)
        \item Standard tournament selection
        \item Evaluated against current state only
    \end{itemize}
\end{columns}
\vspace{0.2cm}\centering
All four evaluated on the \textbf{same dynamically evolving terrain} each generation.
\end{frame}
% FIXED: 30 → 22 gens, 32×32 → 48×48, removed "shifting resources"
\begin{frame}{Coevolution Results --- \textcolor{progAmber}{Pending}}
\small
\textbf{30-Generation Arms Race on Dynamic 48$\times$48 Map:}
\begin{itemize}
    \item Environment: progressive blockades (6--14 walls/gen), mixed wall types
    \item All 4 algorithms co-evolve --- each against the other three's best agent
    \item \textbf{Results:} \textcolor{progAmber}{\textbf{To be collected}} --- coevolution experiments are currently running
\end{itemize}
\vspace{0.2cm}
\textbf{Expected Evaluation:}
\begin{itemize}
    \item Final showdown: 10 games with alternating sides for fairness
    \item Compare evolved strategy profiles (economy vs military weights)
    \item Measure adaptation speed across terrain shifts
\end{itemize}
\vspace{0.2cm}
\textbf{Hypothesis:}
\begin{itemize}
    \item TD-CCEA's trajectory memory should provide an advantage over Standard CCEA's memoryless search in changing terrain
    \item Standard CCEA may memorize one strategy but fail when the map changes
\end{itemize}
\end{frame}
% ──────────────────────────────────────────────────────────────
%  SECTION: Evaluation — Head-to-Head Round-Robin
% ──────────────────────────────────────────────────────────────
\section{Evaluation: Head-to-Head Round-Robin}
% FIXED: 30 → 22 gens, 32×32 → 48×48
\begin{frame}{Evaluation Methodology --- Head-to-Head Tournament}
\small
\begin{block}{How We Measure: Coevolutionary Showdown}
Instead of testing against fixed bots, algorithms are evaluated directly against each other. We run a \textbf{head-to-head tournament} to see which evolutionary method produces the most robust strategies.
\end{block}
\vspace{0.15cm}
\textbf{Two-Phase Evaluation:}
\begin{enumerate}
    \item \textbf{Phase 1 --- Coevolution:} All four algorithms (CCE-TDCCEA, TD-CCEA, SA-TDCCEA, Standard CCEA) co-evolve against each other's best agents for 30 generations on dynamically changing 48$\times$48 maps.
    \item \textbf{Phase 2 --- Final Showdown:} The ultimate evolved agent from each algorithm is pitted against the others in a \textbf{round-robin tournament} (10 head-to-head games per matchup, alternating sides).
\end{enumerate}
\vspace{0.15cm}
\textbf{Fairness Guarantees:}
\begin{itemize}
    \item Sides alternated every game (Player 0 / Player 1 swap)
    \item Server handles direct \texttt{EVOLVED:w1,EVOLVED:w2} matchups natively
    \item Same terrain and resource constraints applied equally to both sides
\end{itemize}
\end{frame}


\begin{frame}{Ablation Study --- What Each Mechanism Contributes}
\end{frame}

\begin{frame}{Wilcoxon test}
\end{frame}

% ════════════════════════════════════════════════════════════
\section{Journal Submission}
% ════════════════════════════════════════════════════════════

\begin{frame}{Journal Submission --- Swarm and Evolutionary Computation}
\begin{block}{Paper Submitted to Elsevier SWEVO (Q1, IF: 8.2)}
Our research paper is currently \textbf{under review} at a top-tier venue for evolutionary and swarm intelligence research.
\end{block}
\vspace{0.2cm}\centering
\renewcommand{\arraystretch}{1.3}
\begin{tabular}{p{4.2cm}p{7.8cm}}
    \toprule
    \textbf{Field} & \textbf{Details}\\
    \midrule
    \textbf{Journal}       & Swarm and Evolutionary Computation (Elsevier)\\
    \textbf{Manuscript ID}  & SWEVO-D-26-01838\\
    \textbf{Title}          & Trajectory-Driven Competitive Co-Evolutionary Algorithm for Strategy Optimization in Competitive Game Environments\\
    \textbf{Date Submitted} & July 7, 2026\\
    \textbf{Status Date}    & July 16, 2026\\
    \textbf{Current Status} & \textcolor{progAmber}{\textbf{Under Review}}\\
    \bottomrule
\end{tabular}
\end{frame}

% ════════════════════════════════════════════════════════════
\section{Implementation Status \& Phase 3 Plan}
% ════════════════════════════════════════════════════════════

\begin{frame}{Implementation Progress}
\scriptsize
\renewcommand{\arraystretch}{1.05}
\begin{tabular}{p{5.2cm}p{2.3cm}p{4.5cm}}
    \toprule
    \textbf{Module} & \textbf{Status} & \textbf{Notes}\\
    \midrule
    Core TD-CCEA algorithm          & \textcolor{doneGreen}{\textbf{Done}} & Trajectory scoring, adaptive mutation\\
    SA-TDCCEA variant               & \textcolor{doneGreen}{\textbf{Done}} & Per-gene sensitivity profiling\\
    CCE-TDCCEA variant              & \textcolor{doneGreen}{\textbf{Done}} & 8D+12D cooperative decomposition\\
    10-algorithm benchmark suite    & \textcolor{doneGreen}{\textbf{Done}} & CMA-ES, DE, PSO, HC, GESY, SA-PDCoEA\\
    Dynamic environment coevolution & \textcolor{doneGreen}{\textbf{Done}} & Progressive blockades + BFS\\
    14-strategy discrete evolution  & \textcolor{doneGreen}{\textbf{Done}} & 32$\times$32 maps\\
    Statistical validation          & \textcolor{doneGreen}{\textbf{Done}} & Wilcoxon signed-rank, 17+ runs\\
    Ablation study                  & \textcolor{doneGreen}{\textbf{Done}} & 15 runs, trajectory + decomposition\\
    Head-to-head round-robin        & \textcolor{progAmber}{\textbf{In Progress}} & 8 algorithms, 280 games\\
    Journal paper submission        & \textcolor{doneGreen}{\textbf{Done}} & SWEVO-D-26-01838, Under Review\\
    StarCraft II bot framework      & \textcolor{progAmber}{\textbf{In Progress}} & Terran, Zerg, Protoss bots\\
    
    \bottomrule
\end{tabular}
\end{frame}


% ════════════════════════════════════════════════════════════
\section{Conclusion}
% ════════════════════════════════════════════════════════════

\begin{frame}{Summary}
\begin{itemize}
    \item \textbf{Phase 1:} Core TD-CCEA with trajectory scoring, adaptive mutation, diversity preservation --- 100\% win rate on 8$\times$8 and 16$\times$16 maps
    \item \textbf{Phase 2:}
    \begin{itemize}
        \item 3 algorithm variants (TD-CCEA, SA-TDCCEA, CCE-TDCCEA)
        \item Dynamic environment coevolution with progressive blockades on 32$\times$32 maps
       
    
        \item \textbf{Journal paper submitted} to SWEVO (Under Review)
    \end{itemize}

\end{itemize}
\end{frame}


% --- REFERENCES ---
\section{References}
\begin{frame}[allowframebreaks]{References}
    \bibliographystyle{IEEEtran}
    \bibliography{ref} 
\end{frame}

% --- THANK YOU ---
\begin{frame}
    \centering    
    {\Huge \textbf{\textcolor{amritaMaroon}{Thank You}}}
\end{frame}


\end{document}